%%
% 1-wegintegral.tex -- 1. Wegintegral
%
% (c) 2025 Prof Dr Andreas Müller
%
\section{Wegintegrale in $\mathbb{C}$
\label{buch:integration:section:wegintegral}}
\kopfrechts{Wegintegrale in $\mathbb{C}$}
Das Riemann-Integral einer Funktion ist definiert für Invervalle.
\index{Riemann-Integral}%
Mit Hilfe eines Weges im Definitionsgebiet einer komplexen Funktion
können wir die Funktion in eine komplexwertige Funktion auf einem
Intervall, nämlich dem Definitionsbereich des Weges verwandeln.
Da es zwischen zwei Punkten von $U$ beliebig viele verschiedene
Wege gibt, müssen die Bedingungen identifiziert werden, unter denen
das Integral nicht von der Wahl des Weges abhängt.

%
% 11-wege.tex -- Wege in $\mathbb{C}$
%
% (c) 2026 Prof Dr Andreas Müller
%
\subsection{Wege in $\mathbb{R}^n$ und $\mathbb{C}$}
Ein \emph{Weg} in einer offenen Menge $U\subset \mathbb{R}^n$ ist eine
\index{Weg}
stetige Abbildung
\[
\gamma : I \to U : t \mapsto \gamma(t)
\]
für ein Intervall $I=[a,b]$.
Der Weg heisst \emph{geschlossen}, wenn Anfangs- und Endpunkt übereinstimmen,
\index{Weg!geschlossen}%
\index{geschlossener Weg}%
wenn also $\gamma(a)=\gamma(b)$ gilt.
Ein Weg heisst \emph{differenzierbarer Weg}, wenn die Funktion $\gamma$
\index{Weg!differenzierbar}
differenzierbar ist.
Die Ableitung
\[
\dot{\gamma}
\colon
I
\to
\mathbb{R}^n
:
t
\mapsto
\dot{\gamma}(t)
=
\frac{d\gamma}{dt}(t)
\]
heisst auch der Tangentialvektor des Weges im Punkt $\gamma(t)$.

%
% Parametertransformation auf Standardintervalle
%
\subsubsection{Parametertransformation auf Standardintervalle}
Manchmal ist es nützlich, spezielle Parameterintervalle zu verwenden.
Die Parametertransformation
\begin{equation}
t
\colon
[0,1] \to [a,b]
:
s
\mapsto
a + s (b-a)
=
(1-s)a
+
sb
\label{buch:integration:wegintegral:eqn:transformation}
\end{equation}
transformiert den Weg $\gamma\colon[a,b]\to \mathbb{R}^n$
auf den Weg $\gamma\circ t\colon [0,1]\to\mathbb{R}^n$.
Der Tangentialvektor ist nach der Kettenregel
\[
\frac{d}{ds}(\gamma\circ t)(s)
=
\dot{\gamma}(t(s))\cdot \frac{dt}{ds}
=
\dot{\gamma}(t(s))\cdot (b-a).
\]
Er wird also mit dem gleichen Faktor $b-a$ skaliert wie der Parameter.

Mit der Transformation
\eqref{buch:integration:wegintegral:eqn:transformation}
kann jeder Wege auf das Standardintervall $[0,1]$ transformiert werden.
Durch Zusammensetzung solcher Transformationen und ihrere Inversen
kann als ein Weg mit jedem beliebigen Parameterintervall
parametrisiert werden.
Wir können daher im Folgenden jederzeit annehmen, dass das Parameterintervall
das Einheitsintervall oder auch jedes beliebige andere Intervall ist.
Wir werden daher meistens das Intervall nicht im Detail spezifizieren
und es dem Leser überlassen, das Intervall aus dem Kontext abzuleiten,
sofern die Wahl des Intervalls überhaupt von Bedeutung ist.

%
% Differenzierbare Wege in C^2
%
\subsubsection{Differenzierbare Wege in $\mathbb{C}^2$}
Da $\mathbb{C}$ als zweidimensionaler reeller Vektorraum betrachtet
werden kann, ist ein Weg in $U\subset\mathbb{C}$ eine Abbildung
von einem Intervall $I$ nach $U$:
\[
\gamma
\colon
I
\to
U
:
t
\mapsto
\gamma(t).
\]
Die Ableitung nach $t$ ist ebenfalls eine komplexe Zahl, die Funktion
\[
\dot{\gamma}
\colon
I
\to
\mathbb{C}
:
t
\mapsto
\dot{\gamma}(t)
=
\frac{d\gamma}{dt}(t)
\]
heisst wieder Tangentialvektor, er ist jetzt aber auch eine komplexe 
Zahl.

%
% Real- und Imaginärteil des Weges
%
\subsubsection{Real- und Imaginärteil des Weges}
Ein Weg in $\mathbb{R}^2$ entsteht dadurch, dass man den Vektor aus
Real- und Imaginärteil von $\gamma(t)$ bildet:
\[
\gamma_{\mathbb{R}}
\colon
I
\to
\mathbb{R}^2
:
t
\mapsto
\begin{pmatrix}
\operatorname{Re}\gamma(t)\\
\operatorname{Im}\gamma(t)
\end{pmatrix}.
\]
Der Tangentialvektor besteht aus den Ableitungen von Real-
und Imaginärteil
\[
\dot{\gamma}
\colon
I
\to
\mathbb{C}
:
t
\mapsto
\frac{d}{dt}\bigl(\operatorname{Re}\gamma(t)\bigr)
+
i
\frac{d}{dt}\bigl(\operatorname{Im}\gamma(t)\bigr)
\]
nach $t$.

\input{chapters/030-integration/fig/fig-wege.tex}%

\begin{beispiel}
\label{buch:integration:wege:bsp:strecke}
Sei $U=\mathbb{C}$ und seien zwei Punkte $z_0,z_1\in \mathbb{C}$ gegeben.
Dann ist
\[
\gamma
\colon
[0,1]
\to
\mathbb{C}
:
(1-t)z_0 + tz_1
\]
ein differenzierbarer Weg, der $\gamma(0)=z_0$ mit $\gamma(1)=z_1$
verbindet.
Die Ableitung von $\gamma$ ist
\[
\dot{\gamma}(t)
=
z_1-z_0.
\]
Insbesondere ist die Ableitung konstant.

Weitere Wege mit den gleichen Endpunkten können konstruiert werden,
indem ein Weg addiert wird, der von $0$ nach $0$ führt.
Zum Beispiel ist für jedes $n\in\mathbb{N}$ 
\[
\delta(t)
=
r_0
(e^{2\pi int} - 1)
\]
ein kreisförmiger Weg mit Radius $r_0$, der $n$ Mal durchlaufen wird.
Fügt man diesen Weg und zusätzlich den Summanden $i\sin \pi t$ zum Weg
$\gamma(t)$ hinzu, erhält man den Weg
\[
\gamma_1(t)
=
\gamma(t)
+
r_0
(e^{2\pi int} - 1)
+
i\sin \pi t,
\]
der in Abbildung~\ref{buch:integration:wege:fig:wege}~a)
{\color{darkgreen}grün} eingezeichnet ist.
Der zugehörige Geschwindigkeitsvektor ist
\[
\dot{\gamma}_1(t)
=
(z_1-z_0)
+
2\pi i
n
r_0
e^{2\pi int}
+
i\pi\cos \pi t.
\qedhere
\]
\end{beispiel}

\begin{beispiel}
\label{buch:integration:wege:bsp:kreis}
Sei $U=\mathbb{C}$ und $r\in\mathbb{R}$.
Der Weg
\[
\gamma
\colon
[0,2\pi]
\to
U
:
t
\mapsto
r(\cos t + i\sin t)
\]
ist ein differenzierbarer Weg mit $\gamma(0) = \gamma(2\pi) = 1$.
Die Ableitung ist
\[
\dot{\gamma}(t)
=
r(-\sin t + i\cos t).
\]
In Satz~\ref{buch:holomorph:holomorph:satz:eulerformel}
wurde gezeigt, dass $\gamma(t) = re^{it}$ ist mit der Ableitung
$\dot{\gamma}(t)=rie^{it}$.
\end{beispiel}

%
% 12-integration.tex -- Integral entlang eines Weges
%
% (c) 2026 Prof Dr Andreas Müller
%
\subsection{Integral entlang eines Weges}
Für eine stetige komplexe Funktion $f(z)$ und einen Weg $\gamma(t)$
ist $f(\gamma(t))$ eine stetige Funktion der reellen Variablen $t$.
Nach allgemeinen Sätzen über das Riemann-Integral können wir also 
\index{Riemann-Integral}%
von Real- und Imaginärteil das Integral über das Definitionsgebiet
berechnen.
Die Definition eines Wegintegrals sollte aber möglichst nicht von
der Parametrisierung des Weges abhängen.
Dies ist nicht selbstverständlich.
Das Integral der konstanten Funktion $1$ z.~B.~ist immer
\begin{equation}
\int_a^b f(\gamma(t)) \,dt
=
\int_a^b 1\,dt
=
b-a,
\label{buch:integration:wegintegral:eqn:konstant}
\end{equation}
hängt also sehr wohl von der Parametrisierung ab.
Daher werden wir zunächst eine von der Parametrisierung unabhängige
Form des Wegintegrals definieren müssen.
Anschliessend werden wir zeigen, wie dieses Wegintegral mit der
Stammfunktion zusammenhängt.
\index{Stammfunktion}

%
% Definition des Wegintegrals
%
\subsubsection{Definition des Wegintegrals}
Das einführende Beispiel
hat gezeigt, dass das Integral
\eqref{buch:integration:wegintegral:eqn:konstant}
von $f(\gamma(t))$ von der Parametrisierung
abhängt.
Bewegt sich der Punkt $\gamma(t)$ langsamer durch das Gebiet $U$,
wird ein längeres Parameterintervall benötigt, um die Strecke
zwischen zwei Punkten zu durchlaufen.
Das gesuchte Wegintegral benötigt also auch noch einen
von der Geschwindigkeit $\dot{\gamma}(t)$ abhängigen Faktor,
um diesen Unterschied auszugleichen.

\begin{definition}[Wegintegral in $\mathbb{C}$]
\label{buch:integration:definition:definition:wegintegral}
Sei $f\colon U\to\mathbb{C}$ eine stetige komplexe Funktion und
$\gamma:I\to U$ ein differenzierbarer Weg $U$ mit $I=[a,b]$.
Dann ist das \emph{Integral von $f$ entlang des Weges $\gamma$}
\index{Integral entlang eines Weges}%
\index{Wegintegral}%
\index{Sgamma@$\int_\gamma$}%
gegeben durch
\begin{equation}
\int_\gamma f(z)\,dz
=
\int_I f(\gamma(t)) \dot{\gamma}(t)\,dt.
\end{equation}
Die gleiche Definition kann verwendet werden, wenn der Weg
$\gamma$ in endlich vielen Punkten nicht differenzierbar ist.
Wenn $\gamma$ ein geschlossener Weg ist, wird das Integral auch
\[
\int_\gamma f(z)\,dz
=
\oint_\gamma f(z)\,dz
\]
geschrieben.
\index{Sogamma@$\oint_\gamma$}%
\end{definition}

Wir schreiben $\gamma(t) = x(t) + iy(t)$ und die Funktion $f=u+iv$ 
als Summe von Real- und Imaginärteil.
Dann wird das Wegintegral zu
\begin{align*}
\int_\gamma f(z)\,dz
&=
\int_a^b \bigl(u(x(t), y(t)) + i v(x(t), y(t))\bigr)\,
(\dot{x}(t)+ i\dot{y}(t))\,dt
\\
&=
\int_a^b u(x(t),y(t)) \dot{x}(t) -  v(x(t),y(t))\dot{y}(t)\,dt
\\
&\qquad
+
i
\int_a^b u(x(t),y(t))\dot{y}(t) + v(x(t),y(t))\dot{x}(t)\,dt,
\end{align*}
und kann dank dieser Formeln allein mit reellen Integralen berechnet
werden.

%
% Unabhängigkeit vond er Parametrisierung
%
\subsubsection{Unabhängigkeit von der Parametrisierung}
Wir müssen überprüfen, dass die
Definition~\ref{buch:integration:definition:definition:wegintegral}
tatsächlich nicht von der Parametrisierung abhängt.
Dazu sei $t(\tau)$ eine invertierbare Funktion, die auf einem Intervall
$I_1=[c,d]$ definiert ist.
Nach der Formel für den Koordinatenwechsel in einem Integral
folgt dann
\begin{align*}
\int_c^d f(\gamma(t(\tau))) \frac{d\gamma(t(\tau))}{d\tau} \,d\tau
&=
\int_c^d f(\gamma(t(\tau))) \dot{\gamma}(t(\tau)) \frac{dt(\tau)}{dt}\,d\tau
\\
&=
\int_{\tau(c)}^{\tau(d)} f(\gamma(t)) \dot{\gamma}(t) \,dt
\\
&=
\int_a^b f(\gamma(t)) \dot{\gamma}(t) \,dt.
\end{align*}
Die Umparametrisierung $\tau(t)$ ändert also den Wert des Integrals
\index{Umparametrisierung}%
nicht.
Insbesonderen hängt der Wert des Integrals nicht vom Parameterintervall ab.
Für die Berechnung des Wegintegrals kann jedes beliebige geeignete
Intervall verwendet werden.
Die Notation $\smash{\int_\gamma}$ hat daher auch keinen Platz, die Wahl
des Intervals anzugeben.

\begin{beispiel}
\label{buch:integration:wege:bsp:kreisintegral}
Wir berechnen das Wegintegral für den geschlossenen Weg von
Beispiel~\ref{buch:integration:wege:bsp:kreis}
für die Funktion $f(z) = z^n$.
Einsetzen der Parametrisierung ergibt
\begin{align*}
\oint_\gamma f(z)\,dz
&=
\int_0^{2\pi}
\bigl( e^{it} \bigr)^{n} ie^{it}
\,dt
=
i
\int_0^{2\pi}
e^{i(n+1)t}
\,dt
=
i
\int_0^{2\pi} \cos (n+1)t + i\sin(n+1)t\,dt.
\end{align*}
Falls $n+1\ne 0$ ist, stehen auf der rechten Seite Integrale der
Sinus- und Kosinusfunktion über eine ganze Anzahl von Perioden, diese
Integrale verschwinden.
Für $n=-1$ verschwindet der Sinus-Term und es bleibt
\[
\oint_\gamma \frac{dz}z
=
i
\int_0^{2\pi}
1\,dt
=
2\pi i.
\]
Zusammengefasst erhalten wir für ganzzahliges $n$
\[
\oint_\gamma z^n\,dz
=
\begin{cases}
0&\qquad\text{falls $n\ne -1$}\\
2\pi i&\qquad\text{falls $n=-1$.}
\end{cases}
\]
Alternativ könnte man das Integral auch mit dem Standardintervall
und dem Weg $\gamma\colon [0,1]\to\mathbb{C}:t\mapsto e^{2\pi it}$
berechnen.
Die Berechnungsformel liefert dann
\begin{align*}
\int_\gamma z^n\,dz
&=
\int_0^1 e^{2\pi i n t}\cdot 2\pi i e^{2\pi i t}\,dt
=
2\pi i
\int_0^1 e^{2\pi i (n+1)t}\,dt
\\
&=
\left\{
\renewcommand{\arraycolsep}{1.5pt}
\begin{array}{rclcllcl}
\displaystyle
2\pi i
\biggl[\frac{1}{2\pi i (n+1)} e^{2\pi i(n+1)t}\biggr]_0^1
&=&
0
&\quad&\text{falls}&
n &\ne& -1,
\\[10pt]
\displaystyle 2\pi i \int_0^1 1\,dt
&=&
2\pi i
&\quad&\text{falls}&
n &=& -1.
\end{array}
\right.
\end{align*}
Die Parametrisierung mit dem Einheitsintervall liefert tatsächlich den
gleichen Wert für das Wegintegral.
\end{beispiel}

%
% Zusammenhängen von Wegen
%
\subsubsection{Zusammenhängen von Wegen}
Seien $\gamma_1$ und $\gamma_2$ mit dem Einheitsintervall parametrisierte
Wege in $U\subset\mathbb{C}$ so, dass der Endpunkt von $\gamma_1$ der
Anfangspunkt von $\gamma_2$ ist, also
\[
\gamma_1(1)
=
\gamma_2(0).
\]
Dann kann man die Verkettung
\index{Verkettung}%
$\gamma_1\cdot\gamma_2$ der beiden Wege durch
\[
\gamma_1\cdot\gamma_2
\colon
I
\to
U
:
t
\mapsto
\begin{cases}
\gamma_1(2t)   &\qquad\text{für $0 \le t \le \frac12$}\\
\gamma_2(2t-1) &\qquad\text{für $\frac12 \le t \le 1$}
\end{cases}
\]
definieren, die ebenfalls durch das Einheitsintervall parametrisiert ist.
Der Weg $\gamma_1\cdot\gamma_2$ führt vom Punkt $\gamma_1(0)$
über $\gamma_1(1)=\gamma_2(0)$ zum Punkt $\gamma_2$.

%
% Umkehrung der Durchlaufrichtung
%
\subsubsection{Umkehrung der Durchlaufrichtung}
Die Umkehrung der Durchlaufrichtung eines Weges hat Auswirkungen
\index{Durchlaufrichtung}%
auf den Wert des Wegintegrals, die wir später brauchen werden.

\begin{definition}[umgekehrter Weg]
\label{buch:integration:wegintegral:def:wegumkehrung}
Sei $\gamma\colon I\to U$ ein mit dem Einheitsintervall parametrisierter
differenzierbarer Weg.
Dann ist
\[
\gamma^{-1}
\colon
I
\to
U
:
t\mapsto \gamma(1-t)
\]
ein differenzierbarer Weg in $U$.
Er heisst der \emph{umgekehrte Weg} von $\gamma$, er wird in
\index{Weg!umgekehrt}%
entgegengesetzer Richtung durchlaufen.
\end{definition}

Eine alternative Möglichkeit, den umgekehrten Weg zu
$\gamma\colon [a,b]\to\mathbb{C}$ zu definieren, ist
\[
\gamma^{-1}\colon [-b,-a] \to \mathbb{C} : t\mapsto \gamma(-t).
\]
Für das Einheitsintervall kann diese Parametrisierung durch eine
Verschiebung des Parameters auf die
Definition~\ref{buch:integration:wegintegral:def:wegumkehrung}
zurückgeführt werden.

Man beachte, dass die Notation $\gamma^{-1}$ auch als die inverse
\index{gamma-1@$\gamma^{-1}$}%
Funktion gelesen werden können, dies ist jedoch nicht, was hier
gemeint ist.
Tatsächlich hat die Umkehrfunktion von $\gamma$ in diesem Zusammenhang
meistens keine sinnvolle Anwendung hat, so dass nicht wirklich
eine Gelegenheit für eine Verwechselung entsteht.

\begin{satz}
\label{buch:integration:wegintegral:satz:wegumkehrung}
Ist $\gamma$ ein differenzierbarer Weg in $U\subset\mathbb{C}$ und
$f\colon U\to\mathbb{C}$ eine stetige Funktion, dann ist
\[
\int_{\gamma^{-1}} f(z)\,dz
=
-
\int_{\gamma} f(z)\,dz.
\]
\end{satz}

\begin{proof}
Da es nicht auf die Parametrisierung des Weges ankommt, gehen wir
von der Parametrisierung durch ein Einheitsintervall aus.
Die Definition~\ref{buch:integration:wegintegral:def:wegumkehrung}
von $\gamma^{-1}$ bedeutet, dass die Ableitung von 
$\gamma^{-1}$
\[
\frac{d}{dt}
\gamma^{-1}(t)
=
\frac{d}{dt}
\bigl(\gamma(1-t)\bigr)
=
\dot{\gamma}(1-t)\cdot(-1)
\]
ist.
Jetzt kann man das Wegintegral
\begin{align*}
\int_{\gamma^{-1}} f(z)\,dz
&=
-
\int_0^1
f(\gamma(1-t))
\dot{\gamma}(1-t)
\,dt
\intertext{mit der Variablentransformation $s=1-t$ als}
&=
-
\int_1^0
f(\gamma(s))
\dot{\gamma}(s)
\,(-ds)
=
-\int_0^1
f(\gamma(s))
\dot{\gamma}(s)
\,ds
\end{align*}
schreiben.
Damit ist der Satz bewiesen.
\end{proof}

%
% Wegintegrale über achsenparallele Wege
%
\subsubsection{Wegintegrale über achsenparallele Wege}
\input{chapters/030-integration/fig/fig-rechteck.tex}%
Wir betrachten eine stetige komplex Funktion $f\colon U\to\mathbb{C}$
und zwei Punkte $a,b\in U$.
Der Einfachheit halber nehmen wir an, dass das Rechteck $R$ mit Ecken $a$
und $b$ in $U$ enthalten ist.
Als Weg zwischen $a$ und $b$ können Teile des Randes des Rechtecks $R$
verwendet werden.

Der Weg $\gamma_1$ in Abbildung~\ref{buch:integration:wegintegral:fig:rechteck}
verläuft zunächst parallel zur reellen Achse, bevor er parallel zur
imaginären Achse in den Punkt $b$ führt.
Der horizontale Teil des Weges kann mit
\begin{align*}
\gamma_{1,h}
&\colon
[\operatorname{Re}a,\operatorname{Re}b]
\to
\mathbb{C}
:
t
\mapsto
t + i\operatorname{Im} a
&&\Rightarrow&
\dot{\gamma}_{1,h} &= 1
\intertext{parametrisiert werden, während der vertikale Teil
die Parametrisierung}
\gamma_{1,v}
&\colon
[\operatorname{Im}a,\operatorname{Im}b]
\to
\mathbb{C}
:
s
\mapsto
\operatorname{Re}b + is
&&\Rightarrow&
\dot{\gamma}_{1,h} &= i
\end{align*}
hat.
Das Wegintegral der Funktion über den Weg $\gamma_1$ kann daher als
Summe zweier Integrale
\begin{align}
F_1(b)
&=
\int_{\gamma_1} f(z)\,dz
\notag
\\
&=
\int_{\gamma_{1,h}} f(z)\,dz
+
\int_{\gamma_{1,v}} f(z)\,dz
\notag
\\
&=
\int_{\operatorname{Re}a}^{\operatorname{Re}b}
f(t+i\operatorname{Im}a)
\,dt
+
\int_{\operatorname{Im}a}^{\operatorname{Im}b}
f(\operatorname{Re}b+is)
\cdot
i
\,ds
\label{buch:integration:wegintegral:achsparallel:eqn:F1}
\end{align}
geschrieben werden.
Die Ableitung von $F_1(b+iy)$ nach $y$ kann damit ebenfalls berechnet
werden:
\begin{align*}
\frac{1}{i}
\frac{\partial}{\partial y} F_1(b+iy)
&=
\frac{1}{i}
\frac{\partial}{\partial y}
\int_{\operatorname{Im}a}^{\operatorname{Im}b+y}
f(\operatorname{Re}b+is)\cdot i
\,ds
=
\frac{\partial}{\partial y}
\int_{\operatorname{Im}a}^{\operatorname{Im}b+y}
f(\operatorname{Re}b+is)
\,ds
\\
&=
f(\operatorname{Re}+i(\operatorname{Im}b+y))
=
f(b).
\end{align*}
Da der Weg $\gamma_1$ im Punkt $b$ vertikal verläuft, kann man ihn
zum Zweck der Ableitung noch etwas über $\operatorname{Im}b$ hinaus
verlängern.
Er definiert so eine Funktion $F_1(b)$, deren
Ableitung nach $i\operatorname{Im}b$ mit $f(b)$ übereinstimmt.

Andererseits erlaubt der Weg $\gamma_2$ eine Funktion
\begin{align}
F_2(b)
&=
\int_{\gamma_2} f(z) \,dz
\notag
\\
&=
\int_{\gamma_{2,v}}
f(z)\,dz
+
\int_{\gamma_{2,h}}
f(z)\,dz
\notag
\\
&=
\int_{\operatorname{Im}a}^{\operatorname{Im}b}
f(\operatorname{Re}a + is) \cdot i \,ds
+
\int_{\operatorname{Re}a}^{\operatorname{Re}b}
f(t + i\operatorname{Im}b)
\,dt
\label{buch:integration:wegintegral:achsparallel:eqn:F2}
\end{align}
von $b$ zu definieren.
Durch horizontale Verlängerung des Weges $\gamma_2$ über den Punkt $b$
hinaus kann man jetzt auch die Ableitung von $F_2(b+x)$ berechnen:
\begin{align*}
\frac{\partial}{\partial x}F_2(b+x)
&=
\frac{\partial}{\partial x}
\int_{\operatorname{Re}a}^{\operatorname{Re}b+x}
f(t + i\operatorname{Im}b)
\,dt
\\
&=
f(\operatorname{Re}b+x+i\operatorname{Im}b)
=
f(b+x).
\end{align*}
Der Weg $\gamma_2$ definiert also eine Funktion $F_2(b)$, deren
Ableitung nach $\operatorname{Re}b$ mit $f(b)$ übereinstimmt.

Wenn wir zusätzlich zeigen können, dass die beiden Funktionen
$F_1(b)$ und $F_2(b)$ übereinstimmen, dann ist dadurch eine
Stammfunktion von $f$ definiert.

%
% Integral einer holomorphen Funktion über ein Rechteck
%
\subsubsection{Integral einer holomorphen Funktion über ein Rechteck}
\index{Rechteck}%
In diesem Abschnitt nehmen wir zusätzlich an, dass Real- und Imaginärteil
als Funktionen von zwei reellen Variablen differenzierbar sind.
Wir möchten die Voraussetzungen finden, unter denen es nicht darauf
ankommt, welchen der zwei betrachteten Wege wir nehmen.
Wir möchten also untersuchen, wann die Differenz $F = F_1(b) - F_2(b)$
verschwindet.

Wir berechnen die Differenz $F_1(b)-F_2(b)$ aus
\eqref{buch:integration:wegintegral:achsparallel:eqn:F1}
und
\eqref{buch:integration:wegintegral:achsparallel:eqn:F2}
\begin{align*}
F_1(b) - F_2(b)
&=
\int_{\operatorname{Re}a}^{\operatorname{Re}b}
f(t+i\operatorname{Im}a)
\,dt
+
\int_{\operatorname{Im}a}^{\operatorname{Im}b}
f(\operatorname{Re}b+is)
\cdot
i
\,ds
\\
&\qquad
-
\int_{\operatorname{Im}a}^{\operatorname{Im}b}
f(\operatorname{Re}a + is) \cdot i
\,ds
-
\int_{\operatorname{Re}a}^{\operatorname{Re}b}
f(t + i\operatorname{Im}b)
\,dt
\\
&=
-
\int_{\operatorname{Re}a}^{\operatorname{Re}b}
f(t + i\operatorname{Im}b)
-
f(t + i\operatorname{Im}a)
\,dt
\\
&\qquad
+
i
\int_{\operatorname{Im}a}^{\operatorname{Im}b}
f(\operatorname{Re}b+is)
-
f(\operatorname{Re}a+is)
\,ds.
\end{align*}
Die Differenzen im Integranden können als Integral einer Ableitung
geschrieben werden:
\begin{equation}
\begin{aligned}
f(\operatorname{Re}b+is)
-
f(\operatorname{Re}a+is)
&=
\int_{\operatorname{Re}a}^{\operatorname{Re}b}
\frac{\partial f}{\partial x}(x+is)
\,dx
\\
f(t + i\operatorname{Im}b)
-
f(t + i\operatorname{Im}a)
&=
\int_{\operatorname{Im}a}^{\operatorname{Im}b}
\frac{\partial f}{\partial y}(t+iy)
\,dy.
\end{aligned}
\label{buch:integration:wegintegral:achsparallel:eqn:partint}
\end{equation}
Damit wird die Differenz
\begin{align*}
F_1(b) - F_2(b)
&=
-
\int_{\operatorname{Re}a}^{\operatorname{Re}b}
\int_{\operatorname{Im}a}^{\operatorname{Im}b}
\frac{\partial f}{\partial y}(x+iy)
\,dy
\,dx
+
i
\int_{\operatorname{Im}a}^{\operatorname{Im}b}
\int_{\operatorname{Re}a}^{\operatorname{Re}b}
\frac{\partial f}{\partial x}(x+iy)
\,dx
\,dy.
\intertext{Da über ein Rechteck integriert wird, dürfen die beiden
Integrationen des zweiten Integrals vertauscht werden, die Differenz
ist also auch}
&=
\int_{\operatorname{Re}a}^{\operatorname{Re}b}
\int_{\operatorname{Im}a}^{\operatorname{Im}b}
-\frac{\partial f}{\partial y}(x+iy)
+
i
\frac{\partial f}{\partial x}(x+iy)
\,dy
\,dx
\\
&=
-2i
\int\hspace*{-2mm}\int_{R}
\frac12
\biggl(
\frac{\partial f}{\partial x}
+
i
\frac{\partial f}{\partial y}
\biggr)
\,dx\,dy
=
2i
\int\hspace*{-2mm}\int_R
\frac{\partial f}{\partial\bar{z}}
\,dx\,dy.
\end{align*}
Die Differenz verschwindet also genau dann, wenn die Ableitung
$\partial f/\partial\bar{z}$ verschwindet, also genau für holomorphe
Funktionen.

\begin{satz}
\label{buch:integration:wegintegral:satz:rechteckintegrale}
Sei $f\colon U\to\mathbb{C}$ 
eine Funktion, deren Real- und Imaginärteile als Funktionen von zwei
reellen Variablen differenzierbar sind. 
Dann sind die Integrale über die Wege $\gamma_1$ und $\gamma_2$ für
beliebige Rechtecke mit den Ecken $a,b\in\mathbb{C}$, die in $U$
enthalten sind, gleich, wenn die Funktion $f$ holomorph ist.
\end{satz}

%
% Dreiseitige Wege
%
\subsubsection{Dreiseitige Wege}
\input{chapters/030-integration/fig/fig-dreiecke.tex}
Der Satz~\ref{buch:integration:wegintegral:satz:rechteckintegrale}
zeigt, dass das Wegintegral einer holomorphen Funktion über den
``linken oberen'' und den ``rechten unteren'' Rand eines Rechtecks 
mit den Ecken $a$ und $b$ gleich ist.
Da sich jede beliebige Kurve durch einen Polygonzug mit achsenparallelen
Segmenten approximieren lässt, ist anzunehmen, dass auch das Integral
über eine beliebige Kurve von $a$ nach $b$ gleich sein dürfte.
Um dies zu zeigen, beginnen wir mit der in
Abbildung~\ref{buch:integration:wegintegral:fig:dreiecke}
dargestellten Situation, dass sich der Weg von $a$ nach $b$ als Graph
der Funktionen $\varphi(x)$ und $\psi(y)$ darstellen lässt.
Solche Gebiete wollen wir im Folgenden \emph{Dreiseite} nennen.
\index{Dreiseit}%
Sowohl
\[
x\mapsto (x,\varphi(x))
\qquad\text{als auch}\qquad
y\mapsto (\psi(y), y)
\]
beschreibt also jeweils die diagonale rote Kurve $\gamma$ in
Abbildung~\ref{buch:integration:wegintegral:fig:dreiecke},
während die achsparallelen Abschnitte mit $\gamma_x$ und $\gamma_y$
bezeichnet werden sollen.

Damit kann man jetzt für jedes der Gebiete in
Abbildung~\ref{buch:integration:wegintegral:fig:dreiecke}
der Unterschied der Wegintegrale über die achsparallelen Teile
und über die gekrümmte Seite durch ein Gebietsintegral 
ausgedrückt werden.
Das Integral über $\gamma$ kann wahlweise durch $x$ im Interval
$[\operatorname{Re}a,\operatorname{Re}b]$ als $x+i\varphi(x)$
oder durch $y$ im Interval
$[\operatorname{Im}a,\operatorname{Im}b]$ als $\psi(y)+iy)$
parametrisiert werden.
Wir verwenden dazu die gleiche Vorgehensweise wie im vorangegangenen
Abschnitt, ersetzen aber den Weg $\gamma_2$ durch den Weg
$\color{darkred}\gamma$, so dass wir das Integral
\begin{align}
F
=
\int_{\gamma} f(z)\,dz
&=
\int_{\operatorname{Re}a}^{\operatorname{Re}b}
f(x+i\varphi(x))\cdot(1+i\varphi'(x))\,dx
=
\int_{\operatorname{Re}a}^{\operatorname{Re}b}
f(\delta_1(x))\cdot\dot{\delta}_1(x)\,dx
\label{buch:integration:wegintegral:eqn:int1}
\\
&=
\int_{\operatorname{Im}a}^{\operatorname{Im}b}
f(\psi(y)+iy)\cdot(\psi'(y)+i)\,dy
=
\int_{\operatorname{Im}a}^{\operatorname{Im}b}
f(\delta_2(y))\cdot\dot{\delta}_2(y)\,dy
\label{buch:integration:wegintegral:eqn:int2}
\end{align}
erhalten, wobei
\[
\begin{aligned}
\delta_1(t) &= t+i\varphi(t)
&&\qquad&
&\text{mit }\phantom{s}\llap{$t$}\in[\operatorname{Re}a,\operatorname{Re}b]
\\
\delta_2(s) &= \psi(s)+is
&&\qquad&
&\text{mit }
s\in[\operatorname{Im}a,\operatorname{Im}b]
\end{aligned}
\]
zwei verschiedene Parametrisierungen der Kurve $\gamma$ sind.

Mit den Funktion $\varphi(t)$ und $\psi(s)$ lassen sich 
Parametertransformationen durchführen, es gilt nämlich
\begin{align*}
\psi(\varphi(t))&=t
&&\text{und}&
\varphi(\psi(s))&=s
\intertext{und daher nach der Kettenregel auch }
\psi'(\varphi(t)) \varphi'(t) &= 1
&&\text{und}&
\varphi'(\psi(s)) \psi'(s) &= 1
\intertext{oder}
\psi'(\varphi(t)) &= \frac{1}{\varphi'(t)}
&&\text{und}&
\varphi'(\psi(s)) &= \frac{1}{\psi'(s)}.
\end{align*}
Mit der Parametertransformation $s=\varphi(t)$ bzw.~$t=\psi(s)$ kann man
die Teile der Integrale mit Ableitungen der Funktionen $\varphi(t)$
und $\psi(s)$ in
\eqref{buch:integration:wegintegral:eqn:int1}
und
\eqref{buch:integration:wegintegral:eqn:int2}
durch
\begin{align*}
\int_{\operatorname{Re}a}^{\operatorname{Re}b}
f(t+i\varphi(t))\, \varphi'(t)
\,dt
&=
\int_{\operatorname{Im}a}^{\operatorname{Im}b}
f(\psi(s)+is)
\,ds
\\
\int_{\operatorname{Im}a}^{\operatorname{Im}b}
f(\psi(s)+is)\,\psi'(s)
\,ds
&=
\int_{\operatorname{Re}a}^{\operatorname{Re}b}
f(t+i\varphi(t))
\,dt
\end{align*}
in solche ohne die Ableitungen von $\varphi(t)$ und $\psi(s)$
transformieren.
Nach der Transformation treten die $\varphi'$- und $\psi'$-Terme in den 
Integranden nicht mehr auf.
Das Integral bekommt jetzt die Form
\begin{align}
F
=
\int_\gamma f(z)\,dz
&=
\int_{\operatorname{Re}a}^{\operatorname{Re}b}
f(t+i\varphi(t))
\,dt
+
i
\int_{\operatorname{Re}a}^{\operatorname{Re}b}
f(t+i\varphi(t))
\,
\varphi'(t))
\,dt
\notag
\\
&=
\int_{\operatorname{Re}a}^{\operatorname{Re}b}
f(t+i\varphi(t))
\,dt
+
i
\int_{\operatorname{Im}a}^{\operatorname{Im}b}
f(\psi(s)+is)
\,ds.
\label{buch:integration:wegintegral:eqn:ableitungslos}
\end{align}

%
% Wegintegrale über die Randstücke von Dreiseiten
%
\subsubsection{Wegintegrale über die Randstücke von Dreiseiten}
Mit der Formel
\eqref{buch:integration:wegintegral:eqn:ableitungslos}
können wir jetzt die Differenz 
\begin{align*}
F_1(b) - F
&=
\int_{\operatorname{Re}a}^{\operatorname{Re}b}
f(t+i\operatorname{Im}a)
\,dt
+
i
\int_{\operatorname{Im}a}^{\operatorname{Im}b}
f(\operatorname{Re}b+is)
\,ds
{\color{darkred}\mathstrut-
F}
\intertext{der Wegintegrale über die blauen und roten Randteile ohne
die Terme mit den Ableitungen $\varphi'$ und $\psi'$ als}
&=
\int_{\operatorname{Re}a}^{\operatorname{Re}b}
f(t+i\operatorname{Im}a)
{\color{darkred}\mathstrut-
f(t+i\varphi(t))}
\,dt
\\
&\hspace*{3.4cm}
+
i
\int_{\operatorname{Im}a}^{\operatorname{Im}b}
f(\operatorname{Re}b+is)
{\color{darkred}\mathstrut-
f(\psi(y)+is)}
\,ds
\intertext{ausdrücken.
Die Differenzen in den Integranden kann man wie in
\eqref{buch:integration:wegintegral:achsparallel:eqn:partint}
als Integral über eine partielle Ableitung ausdrücken.
So entstehen zwei Doppelintegrale}
&=
\int_{\operatorname{Re}a}^{\operatorname{Re}b}
\int_{\operatorname{Im}a}^{\varphi(t)}
-
\frac{\partial f}{\partial y}(t+is)
\,ds
\,dt
+
\int_{\operatorname{Im}a}^{\operatorname{Im}b}
\int_{\psi(y)}^{\operatorname{Re}b}
i
\frac{\partial f}{\partial x}(t+is)
\,dt
\,ds.
\intertext{Die beiden doppelten Integrale erstrecken sich beide
über das Gebiet $G_1$, sie können daher als ein einziges Doppelintegral}
&=
\int_{G_1}
-
\frac{\partial f}{\partial y}(t+is)
+
i
\frac{\partial f}{\partial x}(t+is)
\,dt
\,ds
\intertext{über das Gebiet $G_1$ geschrieben werden.
Der Integrand kann als die Wirtinger-Ableitung von $f$
nach $\bar{z}$ geschrieben werden.
Damit kann man die Differenz}
&=
i
\int\hspace*{-2mm}\int_{G_1}
\frac{\partial f}{\partial x}
+
i
\frac{\partial f}{\partial y}
\,dx\,dy
=
2i
\int\hspace*{-2mm}\int_{G_1}
\frac{\partial f}{\partial\bar{z}}
\,dx\,dy
\end{align*}
als ein zweifaches Integral über das Gebeit $G_1$ ausdrücken.

Eine analoge Rechnung lässt sich auch für die anderen Gebiete
$G_i$, $i=2,\dots,4$ durchführen.
In allen Fällen zeigt sich, dass die Differenz der Integrale
ein Gebietsintegral der Wirtinger-Ableitung $\partial f/\partial\bar{z}$
\index{Gebietsintegral}%
\index{Wirtinger-Ableitung}%
über das Gebiet ist.

\begin{satz}
\label{buch:integration:wegintegral:satz:dreiecke}
Ist $f$ ein holomorphe Funktion $\sigma$ die im Gegenuhrzeigersinn
orientierte Randkurve eines der Gebiete $G_k$ in der 
\index{Randkurve}%
Abbildung~\ref{buch:integration:wegintegral:fig:dreiecke},
dann ist die Differenz der Wegintegrale über die blauen und
roten Teile des Randes jeweils
\begin{equation}
(-1)^k
\int_\gamma f(z)\,dz
=
2i
(-1)^k
\int\hspace*{-2mm}\int_{G_k} 
\frac{\partial f}{\partial\bar{z}}(x+iy)
\,dx\,dy.
\label{buch:integration:wegintegral:eqn:dreiecke}
\end{equation}
\end{satz}

\begin{proof}
Es ist nur noch der Vorzeichenfaktor $(-1)^k$ zu klären.
Die Differenz der beiden Integrale kann als eine Summe von
Wegintegralen über die blauen und roten Teile des Randes
betrachtet wird, wobei der rote Teil in umgekehrter Richtung
durchlaufen wird.
So entsteht ein einziges Wegintegral über den Rand des
Gebiets.

In den Situationen $G_1$ und $G_3$ wird das Gebiet von der
zusammegesetzten Kurve im Gegenuhrzeigersinn umlaufen, während
es in den Situationen $G_2$ und $G_4$ im Uhrzeigersinn umlaufen
wird.
Indem man die Durchlaufrichtung des zusammengesetzten Weges
in der Situation $G_2$ und $G_4$ umkehrt, entsteht auch in diesen
Situationen ein Weg, der das Gebiet im Gegenuhrzeigersinn
umläuft.
Das Vorzeichen in
\eqref{buch:integration:wegintegral:eqn:dreiecke}
kompensiert die Umkehrung der Durchlaufrichtung des Weges.
\end{proof}

%
% Beliebige Wege
%
\subsubsection{Beliebige Wege}
\input{chapters/030-integration/fig/fig-beliebig.tex}
Wir betrachten jetzt ein Gebiet in der komplexen Ebene, welches
von einem geschlossenen Weg $\gamma$ berandet wird, wie es in
Abbildung~\ref{buch:integration:wegintegral:fig:beliebig}
dargestellt ist.
Um das Wegintegral berechnen zu können, muss die Randkurve
$\gamma$ differenzierbar sein.
Damit lässt sich die Kurve aber auch in Segmente unterteilen,
die durch umkehrbare Funktionen $\varphi$ und $\psi$ parametrisiert
werden können.
Jedes solche Segment ist der gekrümmt Rand eines Dreiseits.
Es ist aber möglich, dass die achsparallelen Teile des Randes
des Dreiseits nicht im Gebiet enthalten sind.
Indem das gekrümmte Randstück des Dreiseits weiter unterteilt
wird, lassen sich die Randdreiseite soweit verkleinern, dass 
sich der ganze Rand durch Dreiseite abdecken lassen, die im Gebiet
enthalten sind.

Die Abbildung~\ref{buch:integration:wegintegral:fig:beliebig}
zeigt, wie sich der Teil des Gebietes im Inneren der Randdreiseseite
in Rechtecke zerlegen lässt.
Damit ist es gelungen, das Gebiet in Dreiseite $H_k$ und Rechtecke $R_l$
zu zerlegen.
Das Wegintegral über den Rand jedes dieser Teilstücke kann als ein
zweifaches Integral über das Gebiet
\[
\oint_{\partial H_k} f(z)\,dz
=
\int\hspace*{-2mm}\int_{H_k} \frac{\partial f}{\partial\bar{z}}\,dx\,dy
\quad\text{und}\quad
\oint_{\partial R_l} f(z)\,dz
=
\int\hspace*{-2mm}\int_{R_l} \frac{\partial f}{\partial\bar{z}}\,dx\,dy
\]
ausgedrückt werden.
Die Teilgebietsränder im Inneren des Gebietes werden wie in
Abbildung~\ref{buch:integration:wegintegral:fig:beliebig}
durch die halben Pfeile angedeutet immer zweimal in entgegengesetzter
Richtung durchlaufen.
Nach Satz~\ref{buch:integration:wegintegral:satz:wegumkehrung}
bedeutet dies, dass die die Wegintegrale über
diese Teilstücke entgegengesetzt sind und sich in der Summe wegheben.
Es folgt, dass
\begin{align*}
\oint_{\gamma}f(z)\,dz
&=
\sum_k\oint_{\partial H_k} f(z)\,dz
+
\sum_l\oint_{\partial R_l} f(z)\,dz
\\
&=
\sum_k
\int\hspace*{-2mm}\int_{\partial H_k}
\frac{\partial f}{\partial\bar{z}}
\,dx\,dy
+
\sum_l
\int\hspace*{-2mm}\int_{\partial R_l}
\frac{\partial f}{\partial\bar{z}}
\,dx\,dy
\\
&=
\int\hspace*{-2mm}\int_{G}
\frac{\partial f}{\partial\bar{z}}
\,dx\,dy.
\end{align*}
Damit ist der folgende Satz bewiesen.

\begin{satz}[Cauchy]
\label{buch:integration:wegintegral:satz:cauchy}
\index{Satz!Cauchy-Integralsatz}%
\index{Cauchy-Integralsatz}%
Sei $f$ eine komplexe Funktion mit differenzierbarem Real- und
Imaginärteil und sei $\gamma$ ein geschlossener, differenzierbarer
Weg im Definitionsgebiet der Funktion, der das Gebiet $G$ im
Gegenuhrzeigersinn umläuft.
Dann gilt
\begin{equation}
\oint_\gamma f(z)\,dz
=
2i
\int\hspace*{-2mm}\int_G
\frac{\partial f}{\partial\bar{z}}
\,dx\,dy.
\end{equation}
Für eine holomorphe Funktion $f$ gilt
\[
\oint f(z)\,dz
=
0.
\]
\end{satz}


%
% Der Satz von Green
%
\subsubsection{Der Satz von Green}
\input{chapters/030-integration/fig/fig-green.tex}%
Die Aussage von
Satz~\ref{buch:integration:wegintegral:satz:rechteckintegrale}
ist ein Spezialfall des sehr viel allgemeineren Satzes von Green.

\begin{satz}[Green]
\index{Satz!Green}%
\index{Green, Satz von}%
\label{buch:integration:wegintegral:satz:green}
Sei $G\subset\mathbb{R}^2$ ein Gebiet in $\mathbb{R}^2$, welches
von einer Kurve $\gamma(t)=(x(t),y(t))$ berandet wird.
Seien ausserdem $f$ und $g$ zwei auf einer Umgebung von $G$ definierte,
differenzierbare Funktionen.
Dann gilt
\[
\int_{\gamma} f\,dx + g\,dy
=
\int_I
f(x(t),y(t))\,\dot{x}(t)
+
g(x(t),y(t))\,\dot{y}(t)
\,dt
=
\int_G 
\frac{\partial g}{\partial x}
-
\frac{\partial f}{\partial y}
\,dx\,dy
\]
\end{satz}

\begin{proof}
Wir beweisen die Aussage des Satzes unter der vereinfachenden Annahme,
dass der Rand des Gebietes $G$ durch Funktion $\varphi_{\pm}(x)$ und
$\psi_{\pm}(x)$ wie in
Abbildung~\ref{buch:integration:wegintegral:fig:green}
gezeigt, beschrieben werden kann.

Die Differenz der Funktionswerte von $f$ an den Stellen
$(x,\varphi_+(x))$ und $(x,\varphi_-(x))$ kann als Integral der
partiellen Ableitung nach $y$ entlang der grünen Linie
\begin{align*}
f(x,\varphi_+(x)) - f(x,\varphi_-(x))
&=
\int_{\varphi_-(x)}^{\varphi_+(x)} \frac{\partial f}{\partial y}(x,y)\,dy
\intertext{berechnet werden.
Analog kann man die Differenz}
g(\psi_+(y),y) - g(\psi_-(y),y)
&=
\int_{\psi_-(y)}^{\psi_+(y)} \frac{\partial g}{\partial x}(x,y)\,dx
\end{align*}
schreiben.

Das Wegintegral entlang des Randes ist die Differenz der Integrale
entlang der verschiedenfarbigen Teile des Randes.
Die richtige Orientierung verwendet die Vorzeichen wie folgt:
\begin{equation*}
\renewcommand{\arraycolsep}{2pt}
\begin{array}{rclcl}
\displaystyle
\int_\gamma f\,dx
&=&
\displaystyle
-
\int_{x_-}^{x_+}
f(x,\varphi_+(x))
-
f(x,\varphi_-(x))
\,dx
&=&
\displaystyle
-
\int_{x_0}^{x_1}
\int_{\varphi_-(x)}^{\varphi_+(x)}
\frac{\partial f}{\partial y}(x,y)
\,dy
\,dx
\\[10pt]
\displaystyle
\int_\gamma g\,dy
&=&
\displaystyle
\phantom{-}
\int_{y_-}^{y_+}
g(\psi_+(y),y)
-
g(\psi_-(y),y)
\,dy
&=&
\displaystyle
\phantom{-}
\int_{y_-}^{y_+}
\int_{\psi_-(y)}^{\psi_+(y)}
\frac{\partial g}{\partial x}(x,y)
\,dx
\,dy.
\end{array}
\end{equation*}
Die Summe der beiden Gleichungen ergibt die Aussage des Satzes.
\end{proof}


%
% Wegintegral und Stammfunktion
%
\subsubsection{Wegintegral und Stammfunktion}
Der Satz~\ref{buch:integration:wegintegral:satz:green}
zeigt, dass es nicht darauf ankommt, welchen Weg man zwischen
zwei Punkten $a$ und $b$ wählt, solange zwei solche Wege ein
Teilgebiet des Definitionsgebietes beranden.
In diesem Fall lässt sich also durch Wegintegration immer eine
Stammfunktion finden.
\index{Stammfunktion}%
Dazu muss es aber zwischen zwei beliebigen Punkten von $U$ einen
Weg geben.

\begin{definition}[wegzusammenhängend]
Ein Gebiet $U\subset\Omega$ heisst \emph{wegzusammenhängend}, wenn
\index{wegzusammenhangend@wegzusammenhängend}%
es zu je zwei Punkten $a,b\in U$ einen Weg gibt, der $a$ und $b$
verbindet.
\end{definition}

Ein Stammfunktion kann jetzt konstruiert werden, indem man einen
Weg von $a$ nach $b$ wählt und dann die Stammfunktion als das
Wegintegral
\[
F(b)
=
\int_\gamma f(z)\,dz
\]
definiert.
Wir müssen uns davon überzeugen, dass verschiedene Wege den gleichen
Wert $F(b)$ liefern.

Zwei verschiedene Wege $\gamma_1$ und $\gamma_2$
könnten jedoch zu zwei verschiedenen Werten
\[
F_1(b)
=
\int_{\gamma_1}f(z)\,dz
\qquad\text{und}\qquad
F_2(b)
=
\int_{\gamma_2}f(z)\,dz
\]
der Stammfunktion Anlass geben.
Setzt man den Weg $\gamma_1$ mit dem in umgekehrter Richtung
durchlaufenen Weg $\gamma_2^{-1}$ zusammen, entsteht ein
geschlossener Weg $\gamma=\gamma_1\cdot\gamma_2^{-1}$.
Der Unterschied der beiden Werte
\[
F_1(b) - F_2(b)
=
\int_\gamma f(z)\,dz
\]
kann mit Satz~\ref{buch:integration:wegintegral:satz:cauchy}
berechnet werden, wenn der zusammengesetzte Weg $\gamma$ eine
Teilmenge von $\mathbb{C}$ berandet, in dem $f$ definiert ist.
Ist $f$ ausserdem holomorph auf dieser Teilmenge, verschwindet
die Differenz.
Für die Definition einer Stammfunktion ist also notwendig, dass
zwei Wege von $a$ nach $b$ eine Teilmenge von $\mathbb{C}$
beranden, die im Definitionsgebiet der holomorphen Funktion $f$
enthalten ist.

\begin{satz}[Stammfunktion]
\label{buch:integration:wegintegral:satz:stammfunktion}
\index{Stammfunktion}%
\index{Satz!Stammfunktion}%
Sei $U\subset\mathbb{C}$ ein wegzusammenhängendes Gebiet und $a\in U$
ein Punkt.
Aussserdem habe $U$ die Eigenschaft, dass für jeden Punkt $b\in U$
zwei verschiedene Wege, die $a$ und $b$ verbinden, ein offenes Teilgebiet
von $U$ beranden.
Dann hat eine holomorphe Funktion $f\colon U\to\mathbb{C}$ eine Stammfunktion
\[
F(b)
=
\int_\gamma f(z)\,dz,
\]
wobei $\gamma$ ein beliebiger Weg ist, der $a$ mit $b$ verbindet.
\end{satz}

%
% 13-homotopie.tex -- Homotopie
%
% (c) 2026 Prof Dr Andreas Müller
%
\subsection{Homotopie
\label{buch:integration:wegintegral:subsection:homotopie}}
Im vorangegangenen Abschnitt wurde klar, dass das unter geeigneten
Voraussetzungen an das Gebiet das Wegintegral für verschiedene Wege 
zwischen $a$ und $b$ gleich ist.
Die Bedingung war etwas schwerfällig zu formulieren, ergab sich aber
unmittelbar aus dem Satz von Green.
Es musste verlangt werden, zwei Wege zwischen $a$ und $b$ ein in
$U$ enthaltenes Teilgebiet beranden.

Eine etwas intuitivere Formulierung wäre, dass sich die beiden
Wege ineinander deformieren lassen.
Während der Deformation wird das Teilgebiet überstrichen.
\index{Deformation}%
Wir versuchen jetzt, dieses Idee mit einer mathematisch präzisen
Definition zu erfassen.
\input{chapters/030-integration/fig/fig-homotopie.tex}

\begin{definition}[Homotopie von Wegen]
\label{buch:integration:definition:homotpie}
Eine \emph{Homotopie} von Wegen $\gamma_0\colon I\to U$
\index{Homotpie}%
und
$\gamma_1\colon I\to U$, $I=[0,1]$,
in einem Gebiet $U$
mit gleichen Endpunkten
$\gamma_0(0)=\gamma_1(0)$
und
$\gamma_0(1)=\gamma_1(1)$
ist eine Abbildung
\[
H
\colon
I\times I \to U
:
(s,t) \mapsto H(s,t)
\]
mit $H(0,t)=\gamma_0(t)$ und $H(1,t)=\gamma_1(t)$
und konstanten Endpunkten
$H(s,0)=\gamma_i(0)$
und
$H(s,1)=\gamma_i(1)$.
Zwei Wege $\gamma_i\colon I\to U$  heissen \emph{homotop}, in Zeichen
\index{homotop}
$\gamma_0\simeq \gamma_1$, wenn es eine
Homotopie gibt, die die beiden Wege verbindet.
\end{definition}

\begin{satz}
\label{buch:integration:satz:homotopie}
Sei $f\colon U\to\mathbb{C}$ eine holomorphe Funktion und
\[
H
\colon
I\times I
\to
U
\]
eine Homotopie der Wege $\gamma_i$, $i=0,1$, mit $H(s,0) = a$ und $H(s,1)=b$
für $s\in I$.
Dann ist
\[
\int_{\gamma_0} f(z)\,dz
=
\int_{\gamma_1} f(z)\,dz.
\]
\end{satz}

\begin{proof}
Da die Homotopie $H$ eine Abbildung nach $U$ ist, verlaufen
alle Wege
\[
\gamma_s
\colon
I\to\mathbb{C}
:
t
\mapsto
\gamma_s(t) = H(t,s)
\]
in $U$.
Die Vereinigung aller Wege
\[
\bigcup_{0\le s\le 1} \operatorname{im}\gamma_s
=
\{
H(t,s)
\mid
t,s\in I
\}
\subset U
\]
ist in $U$ enhalten, so dass zwei beleibige Kurven immer eine
offene Teilmenge $U$ beranden.
Der Satz~\ref{buch:integration:wegintegral:satz:stammfunktion}
bedeutet daher, dass alle Wegintegrale entlange $\gamma_s$
übereinstimmen, insbesondere auch
\[
\int_{\gamma_0} f(z) \, dz
=
\int_{\gamma_s} f(z) \, dz
= 
\int_{\gamma_1} f(z) \, dz.
\]
Damit ist der Satz bewiesen.
\end{proof}

Der Satz~\ref{buch:integration:satz:homotopie}
sagt, dass ineinander deformierbare Wege zum gleichen Wegintegral führen.
Die Bedingung ist jedoch nicht immer einfach zu verifizieren.
Einfacher wird es, wenn ein Weg zusammenziehbar ist im Sinne
der folgenden Definition.

\begin{definition}[zusammenziehbar]
\label{buch:integration:wegintegral:definition:zusammenziehbar}
Ein geschlossener Weg $\gamma$ in einem Gebiet $U$ heisst
\emph{zusammenziehbar}, wenn der homotop ist zum konstanten Weg
\index{zusammenziehbar}%
$t\mapsto \gamma_1(t)=\gamma(0)=\gamma(1)$.
\end{definition}

In $\mathbb{C}$ ist jeder Weg zusammenziehbar, indem man die Homotopie
\[
H(t,s)
=
(1-s)
\gamma(0)
+
s
\gamma(t) 
\]
verwendet.
Für $s=0$ ist $H(t,0) = \gamma(0)$ der konstante Weg im Punkt
$\gamma(0)$, währen für $s=1$
$H(t,1)=\gamma(t)$ der Weg $\gamma$ ist.
Ein Weg ist nicht zusammenziehbar, wenn er ein ``Loch'' in
der Menge umläuft, da er über das Loch hinweggezogen werden
müsste.
Zum Beispiel ist die Funktion $f(z)=\frac1z$ holomorph in
der Menge $U=\mathbb{C}\setminus\{0\}$ definiert, die an der Stelle
$0$ ein Loch hat.
Ein Weg, der den Nullpunkt umläuft, kann nicht in einen Punkt
zusammengezogen werden, da er dabei den Nullpunkt überqueren müsste.

Mit dem Begriff des zusammenziehbaren Weges ist jetzt
der Satz~\ref{buch:integration:satz:homotopie}
sehr viel anschaulicher zu formulieren.
\input{chapters/030-integration/fig/fig-deformation.tex}

\begin{satz}[Cauchy]
\index{Satz!Cauchy-Integralsatz}%
\index{Cauchy-Integralsatz}
\label{buch:integration:wegintegral:satz:zusammenziehbar}
Ist $\gamma$ ein zusammenziehbarer Weg im Definitionsgebiet einer
holomorphen Funktion, dann gilt
\[
\oint_\gamma f(z)\, dz = 0.
\]
\end{satz}

\begin{proof}
Sei $H$ ein Homotopie zwischen dem konstanten Weg im Punkt $a$ und
der geschlossenen Kurve $\gamma$, also
\[
H(t,0) = a
\qquad
\text{und}
\qquad
H(t,1) = \gamma(t).
\]
Ausserdem soll jeder Weg $t\mapsto H(t,s)$ ein geschlossener Pfad im
Punkt $a$ sein, also $H(0,s)=H(1,s)=a$.
Abbildung~\ref{buch:integration:wegintegral:fig:deformation}
ist die Kurve $\gamma$ rot eingezeichnet.

Der Punkt $b=H(\frac12,1)$ teilt die Kurve $\gamma$ in zwei Kurven
$\gamma_1$ und $\gamma_2$ von $a$ nach $b$ auf.
Die Abbildung $s\mapsto H(\frac12,s)$ ist ein Weg von $a$ nach
$b=H(\frac12,1)$, der in
Abbildung~\ref{buch:integration:wegintegral:fig:deformation}
grün eingezeichnet ist.
Wir wollen zeigen, dass beide Wege $\gamma_1$ und $\gamma_2$ sich in
die grüne Kurve $s\mapsto H(\frac12,s)$ deformieren lässt.

Die Funktion 
\[
K_1
\colon
I\times I
\to
\mathbb{C}
:
(t,s)
\mapsto
K_1(t,s)
=
\begin{cases}
H(t/2s,s)    &\qquad\text{für $t  <  s$}
\\
H(\frac12, t)&\qquad\text{für $s \ge t$}
\end{cases}
\]
deformiert den Weg $\gamma_1$ in den grünen Weg,
denn für $s=1$ entsteht der Weg $\gamma_1$ und
für $s=0$ der grüne Weg.
Für $s$-Werte zwischen $0$ und $1$ folgt der Weg zuerst
der Kurve $t\mapsto H(t,s)$ bis zur grünen Kurve, dann der grünen
Kurve.

Eine ähnliche Deformation des Weges $\gamma_2$ in die grüne
Kurve wird durch
\[
K_2
\colon
I\times I
\to
\mathbb{C}
:
(t,s)
\mapsto
K_2(t,s)
=
\begin{cases}
H((1-2t)/2s,s)    &\qquad\text{für $t  <   s$}
\\
H(\frac12, t)     &\qquad\text{für $s \ge t$}
\end{cases}
\]
gegeben.

Da beide Wege $\gamma_1$ und $\gamma_2$ homotop sind zur
grünen Kurve, sind auch die Wegintegrale entlang der Kurven $\gamma_1$
und $\gamma_2$ gleich dem Wegintegral entlang der grünen Kurve.
Das Wegintegral entlang der Kurve $\gamma$ verschwindet daher.
\end{proof}



